\section{Lean development and the remaining object proofs}
\label{sec:lean}
The repository uses \texttt{PR k} for $k$-ary primitive recursive function
\emph{syntax}, \texttt{PTm n} for terms with $n$ free numerical variables,
\texttt{Fm PTm n} for formulas, and \texttt{Prf PRAAx} for their
derivations~\cite{proveit}. An assertion proved in Lean about a
\texttt{Field} differs in type from a derivation in this embedded theory.
The archive contains examples of both. The distinction determines which
steps can already be imported.

\subsection{What the included Lean modules prove}
\begin{description}[leftmargin=0pt,labelsep=0pt,labelwidth=0pt,style=nextline]
\item[\texttt{MetricIdentities.lean}.]
\texttt{gramClosed} and \texttt{midpointClosed} are closed embedded
iPRA derivations. Each of four integer coordinates is encoded by a
difference of two natural variables. The first proves
$N(u)N(v)=(u\cdot v)^2+[u,v]^2$. The second proves the parallelogram law
$N(u+v)+N(u-v)=2N(u)+2N(v)$. The universal closures have eight
natural binders. The open versions allow substitution of arbitrary
integer terms in an arbitrary context. These endpoints use the repository's
\texttt{ZT.derive}, \texttt{ipra\_ring} and \texttt{closeUniv}.
This is precisely the polynomial part of the length and frame arguments.
Transferring the identities to scalar records also needs their field laws.
\item[\texttt{Translation.lean}, \texttt{Axioms.lean} and \texttt{Bridge.lean}.]
The first module defines the two formula languages, term evaluation,
environments, shadowing and the quantifier translation. Its
\texttt{translate\_correct} proves preservation and reflection for every
formula. The second module works over a linearly ordered field; it
proves the equivalence of affine betweenness and the determinant/dot
test, inner Pasch, the Euclidean axiom and five segments. The bridge
composes these proofs using an independent affine interpretation.
\texttt{translated\_innerPasch} checks the closed translated sentence
including the existential point binder.
\item[\texttt{FrameRelations.lean} and \texttt{CanonicalFrame.lean}.]
The first defines all point-language construction graphs. Its
\texttt{BiGraph} explicitly carries the four existence/uniqueness clauses
for a relational bijection and gives both round trips and quantifier
transport. The second proves the midpoint and collinearity criteria for
a concrete Cartesian plane over a linearly ordered field, the equivalence
of its betweenness test with affine interpolation, the standard-frame
addition, multiplication and order formulas, and the coordinate bijection.
The arbitrary-frame formulas in Section~\ref{sec:frame-inverse} are written
as mathematical proofs.
\item[\texttt{ObjectLanguage.lean}.]
This compiler takes seven PR symbols (validity, the scalar zero, three
ring operations, equality and order) and sends the translated language
to actual \texttt{Fm PTm n}. Under binders it shifts older terms and
guards the new variable; it interprets both scalar tests by the Boolean
numeral one. \texttt{closeTranslation} yields an exact type of sentence
for a closed point formula once the symbols are filled. The scalar symbols
are an explicit input to this file.
\end{description}
The included Lean files were compiled using Lean and mathlib 4.32.0.
The two closed object proofs print the standard Lean metatheoretic
dependencies \texttt{propext}, \texttt{Classical.choice} and
\texttt{Quot.sound}. Their values inhabit \texttt{Prf PRAAx};
the constructors of that inductive calculus fix the available
object axioms. The general scalar presentation and the elimination
theorem here have human proofs; their complete object-calculus proof
generators remain to be implemented.

\subsection{A route to closing the general theorem}
Here is an implementation order and the statement to verify at each step.
All semantic equalities between two codes in these obligations abbreviate
the numerical Boolean predicate for equality of their values.
\begin{enumerate}
\item \emph{Arithmetic and records.} Use the upstream
\texttt{Core/Pairing}, \texttt{Core/Division}, \texttt{Core/Gcd},
\texttt{Core/Order}, \texttt{Core/Search} and \texttt{Core/Definitions}
for pairs, bounded loops and simultaneous recursion. Pick one list format,
for example the one in Section~\ref{sec:scalar}; use pairs for signed
numerators and positive denominators for rationals. Store a triple as
a list of length three. A normalized polynomial list has its final
coefficient nonzero unless it is empty. Define \texttt{valid}, interval
reading and every guard by the exact rational tests in
\eqref{eq:intervalcheck}. Encode the recipe to read a record as a
single PR function of the record and the precision. Prove nesting,
width and endpoint bounds by a Boolean invariant on the bisection counter.
For a bounded universal check, take the conjunction of the Boolean
values by recursion; prove its elimination and introduction lemmas once.
\item \emph{Scalar certificates.} Implement the finite convolution,
signed polynomial division, extended gcd and finite matrix determinant
with their coefficient-vector invariants. For each provisional operation
in the paper define its own fixed PR reader with parameters.
Prove the positive slope estimate, the sign gap, and the \emph{conversion}
lemma: the returned code is valid and agrees with the input recipe
at every precision. Prove agreement by fixing the precision before
induction and budgeting all intermediate reads. In the determinant
construction verify the adjugate coefficient identity by pairing
permutations; use it to certify the monic annihilator in the normalizer.
Prove validity, extensionality, field laws and order laws for
\texttt{add, sub, mul, eqTest, leTest}; add totalized inverse for use in
algorithms, keeping the ring language of the output compiler unchanged.
\item \emph{Root enumeration.} For a list $f$ of scalar coefficients
compute the tag for an identically zero polynomial or a strictly
increasing list $L_f$ of roots. The target on the finite-list branch is
$\mathrm{Valid}(x)\to(f(x)=0\leftrightarrow
     \bigvee_{i<|L_f|}x=L_f[i])$. Here the variable-size disjunction
is a bounded Boolean fold. The finite algebra and rational grid in
Section~\ref{sec:scalar} supply both inclusions. Polynomial IVT is a
separate witness theorem with two input signs; no appeal to a general
real-analysis library is needed.
\item \emph{One elimination.} Define signed remainders by a
degree-clocked loop. Define \texttt{query}, the sign counts
in~\eqref{eq:counts}, and the two tail signs. Prove~\eqref{eq:query}
by a finite sweep over the ordered roots from the preceding step.
The invariant is the integer equality between variation loss and root
contributions through the current cell. For a specialized polynomial
list, prove equivalence between each attainable sign condition and its
computed positive count or tail flag. Define \texttt{pick} from the actual
ordered point/interval samples to return a concrete satisfying code.
\item \emph{Reification and composition.} Generate the coefficient
decision tree with formal numerator/denominator pairs. At each node the
guard records every denominator needed to that point. For a fixed
matrix construct a \texttt{Prf} equivalence between an eliminated binder
and its disjunction of accepting branch guards, using \texttt{pick} on
a positive existential branch. Use the same witnesses for counterexamples
to universal claims. In the external proof generator recurse on a
particular syntax tree. Discharge its connectives by the natural
deduction constructors and the Boolean decision for its children.
For a closed sentence the final goal is a value of \texttt{Prf PRAAx []}
at the output of \texttt{closeTranslation}, or at its negation.
For parameters, universally close the valid-code guards and the
bi-implication with the decision bit. Only a completed object proof
is eligible for transport using the repository's \texttt{pra\_iff\_native}.
\end{enumerate}

The ProveIt packet \texttt{eight\_geometry\_theorems\_ipra} treats exact
rational/projective coordinate calculations with denominator guards.
Its approach applies to pieces such as the line-intersection formula.
The present development needs real algebraic coefficients for circle
intersections and the full first-order transfer. We reuse the repository's
checked arithmetic interface for the two embedded proofs.
The executable companions cover scalar records, exact constructions, the
one-variable signed-remainder computation, parameter-dependent elimination
for prenex polynomial formulas, and exact witness selection. The
optional symbolic routines use SymPy for rational polynomial arithmetic
and isolation. Their tests exercise the algorithms independently of Lean.

\paragraph{Further questions.}
How small can the proof generator for a single quantified sign test be
made using the repository's existing bounded-sum proofs? Can a
specialized compiler for ruler-and-compass statements avoid general
polynomial elimination while emitting comparable object-calculus
certificates? For analytic names, which operations admit continuous
computable selections through degeneracies after the input carries extra
orientation or separation data?
